% Figure from MATH 591 notes, section 3. Compile with the notes preamble.
\begin{tikzpicture}[scale=1.0]
% upstairs: two copies of R
\draw[figB, thick] (-3,1) -- (3,1); \draw[figB, thick] (-3,0) -- (3,0);
\node[figB, font=\scriptsize, anchor=west] at (3.05,1) {$\mathbb{R}\times\{1\}$};
\node[figB, font=\scriptsize, anchor=west] at (3.05,0) {$\mathbb{R}\times\{2\}$};
% identifications (x,1) ~ (x,2), x != 0
\foreach \x in {-2.8,-2.3,-1.8,-1.3,-0.8} { \draw[black!40, densely dotted] (\x,0) -- (\x,1); }
\foreach \x in {2.4,2.8} { \draw[black!40, densely dotted] (\x,0) -- (\x,1); }
% the sequence x = 1/n (scale 2): glued pairs converging to the origins
\foreach \x in {2,1,0.6667,0.5,0.4,0.3333} { \draw[figB!70, thin] (\x,0) -- (\x,1); \fill[figB] (\x,0) circle (1.1pt); \fill[figB] (\x,1) circle (1.1pt); }
\foreach \x in {0.2857,0.25,0.2222,0.2} { \draw[figB!50, very thin] (\x,0) -- (\x,1); }
\node[figB, font=\scriptsize, anchor=south] at (1.65,1.03) {$(\tfrac1n,1)\sim(\tfrac1n,2)$};
% the two origins: not identified
\fill[figR] (0,1) circle (1.7pt); \fill[figR] (0,0) circle (1.7pt);
\node[figR, font=\scriptsize, anchor=south] at (-0.05,1.03) {$(0,1)$};
\node[figR, font=\scriptsize, anchor=north] at (0,-0.03) {$(0,2)$};
\node[figR, font=\scriptsize, anchor=center] at (-0.4,0.5) {$\not\sim$};
% the quotient map
\draw[-stealth, thin, black!55] (0,-0.5) -- node[right, font=\scriptsize, black] {$\pi$} (0,-1.15);
% downstairs: the line with two origins
\draw[figB, thick] (-3,-1.85) -- (-0.1,-1.85); \draw[figB, thick] (0.1,-1.85) -- (3,-1.85);
\fill[figR] (0,-1.62) circle (1.7pt); \fill[figR] (0,-2.08) circle (1.7pt);
\node[figR, font=\scriptsize, anchor=south] at (0,-1.6) {$0_1$};
\node[figR, font=\scriptsize, anchor=north] at (0,-2.1) {$0_2$};
\node[font=\scriptsize, anchor=west] at (3.05,-1.85) {$X/{\sim}$};
\end{tikzpicture}
